% !PRINTED 668
\weberpage{668}
2. \gesperrt{Unter dem Produkt $\mathfrak{A}\mathfrak{B}$ zweier Polygone versteht man das Polygon, das die Punkte von $\mathfrak{A}$ und $\mathfrak{B}$ zugleich enthält, und zwar einen Punkt $\mathfrak{P}$, der $r$mal in $\mathfrak{A}$ und $s$mal in $\mathfrak{B}$ vorkommt, $r+s$mal. Die Ordnung eines Produktes ist gleich der Summe der Ordnungen der Faktoren.}

Wenn wir also mit $\mathfrak{P}_1, \mathfrak{P}_2, \ldots$ verschiedene Punkte bezeichnen, so können wir jedes Polygon $\mathfrak{A}$ auf eine Weise in die Form setzen:

\weq{1}{\mathfrak{A} = \mathfrak{P}_1^{r_1}\mathfrak{P}_2^{r_2}\cdots}

Die Punkte sind also in der Rechnung mit Polygonen die \gesperrt{Primelemente}. Um auch für die Einheit einen Vertreter zu haben, muß noch das \glqq Nulleck\grqq{} $\mathfrak{O}$, das gar keinen Punkt enthält, mitgenannt werden.

\gesperrt{Der größte gemeinschaftliche Teiler zweier Polygone} $\mathfrak{A}, \mathfrak{B}$ enthält jeden Punkt, der in $\mathfrak{A}$ und $\mathfrak{B}$ vorkommt, und zwar so oft, als er in jedem von beiden vorkommt.

\gesperrt{Das kleinste gemeinschaftliche Vielfache} enthält jeden Punkt von $\mathfrak{A}$ und $\mathfrak{B}$, und zwar in der höchsten der Ordnungen, in denen er in $\mathfrak{A}$ und $\mathfrak{B}$ vorkommt.

3. \gesperrt{Ist $z$ eine Variable des Körpers $\Omega$ und $n$ der Grad des Körpers in bezug auf $z$, so nimmt $z$ jeden Wert $c$ in $n$ Punkten an.}

Denn zerlegt man $z - c$ in seine Primfaktoren [§ 183, (17)]:
\[
z - c = \pi_1^{e_1}\pi_2^{e_2}\cdots, \quad e_1 + e_2 + \cdots = n,
\]
so erzeugt jeder dieser Primfaktoren einen Punkt, z.\,B. $\pi_1$ den Punkt $\mathfrak{P}_1$, in dem die Funktion $z - c$ Null in der $e_1$ten Ordnung wird. Rechnet man diesen Punkt $e_1$fach, so erhält man ein Polygon $\mathfrak{P}_1^{e_1}\mathfrak{P}_2^{e_2}\cdots$ von der Ordnung $n$, in dessen Punkten die Funktion $z - c$ verschwindet.

Derselbe Satz gilt aber auch für den Wert $c = \infty$, wie man daraus schließt, daß der Grad des Körpers $\Omega$ in bezug auf $1 : z$ derselbe ist wie in bezug auf $z$.

4. \gesperrt{Die $n$ Punkte, in denen $z$ einen Wert $c$ annimmt, heißen konjugiert nach $z$, und die Werte $\eta_1, \eta_2, \ldots, \eta_n$, die eine Funktion $\eta$ in $\Omega$ in $n$ konjugierten Punkten annimmt, heißen gleichfalls konjugiert nach $z$, und $z$ heißt von der Ordnung $n$. Die Konstanten, und nur diese, haben die Ordnung 0.}
